\chapter{Probabilités et variables aléatoires}

\noindent\textit{Le calcul des probabilités modélise le hasard et l'incertitude :
tirages, jeux, contrôle de qualité, dépistage médical. Après le dénombrement, on
y associe à chaque issue une mesure de sa « chance » de se produire. La série C
approfondit la \textbf{probabilité conditionnelle}, la \textbf{formule des
probabilités totales} et l'étude des \textbf{variables aléatoires}, jusqu'à la
\textbf{loi binomiale}, modèle d'une répétition d'épreuves indépendantes.}
\vspace{8pt}

% ═══════════════════════════════════════════════════════════════
\section{Vocabulaire et équiprobabilité}

\subsection{Expérience aléatoire, événements}

\begin{definition}[Univers et événements]
Une \textbf{expérience aléatoire} est une expérience dont on ne peut prévoir le
résultat. L'ensemble de toutes ses issues forme l'\textbf{univers} $\Omega$. Un
\textbf{événement} est une partie de $\Omega$. Un \textbf{événement élémentaire}
est un singleton $\{\omega\}$. L'événement $\Omega$ est \textbf{certain},
$\varnothing$ est \textbf{impossible}.
\end{definition}

\begin{definition}[Opérations sur les événements]
Soient $A$ et $B$ deux événements.
\begin{itemize}
\item $A\cap B$ (« $A$ \textbf{et} $B$ ») est réalisé si $A$ et $B$ le sont ;
\item $A\cup B$ (« $A$ \textbf{ou} $B$ ») est réalisé si l'un au moins l'est ;
\item $\overline A$, l'\textbf{événement contraire}, est réalisé quand $A$ ne l'est pas ;
\item $A$ et $B$ sont \textbf{incompatibles} si $A\cap B=\varnothing$.
\end{itemize}
\end{definition}

\begin{definition}[Probabilité]
Une \textbf{probabilité} $P$ sur un univers fini $\Omega$ associe à chaque
événement $A$ un réel $P(A)\in[0,1]$ tel que $P(\Omega)=1$ et, pour deux
événements incompatibles, $P(A\cup B)=P(A)+P(B)$.
\end{definition}

\begin{propriete}[Règles de calcul]
Pour tous événements $A$ et $B$ :
\[
P(\overline A)=1-P(A),\qquad
P(A\cup B)=P(A)+P(B)-P(A\cap B),\qquad P(\varnothing)=0.
\]
\end{propriete}

\subsection{Équiprobabilité}

\begin{definition}[Hypothèse d'équiprobabilité]
Lorsque tous les événements élémentaires ont la même probabilité, on dit qu'il y
a \textbf{équiprobabilité}. Pour tout événement $A$ :
\[
P(A)=\frac{\Card A}{\Card\Omega}
=\frac{\text{nombre de cas favorables}}{\text{nombre de cas possibles}}.
\]
\end{definition}

\begin{remarque}
L'équiprobabilité ramène le calcul d'une probabilité à un \textbf{dénombrement}
(chapitre 12) : on compte les issues favorables et les issues possibles à l'aide
des $p$-listes, arrangements ou combinaisons selon que l'ordre et la répétition
interviennent ou non.
\end{remarque}

\begin{exemple}
On tire au hasard une carte d'un jeu de $32$ cartes. Soit $A$ : « obtenir un
roi ». Il y a $4$ rois, donc $P(A)=\dfrac{4}{32}=\dfrac18$. Soit $B$ : « obtenir
un c\oe ur » ; $P(A\cup B)=P(A)+P(B)-P(A\cap B)=\dfrac{4}{32}+\dfrac{8}{32}-\dfrac{1}{32}=\dfrac{11}{32}$.
\end{exemple}

% ═══════════════════════════════════════════════════════════════
\section{Probabilité conditionnelle et indépendance}

\begin{definition}[Probabilité conditionnelle]
Soit $B$ un événement de probabilité non nulle. La \textbf{probabilité
conditionnelle de $A$ sachant $B$} est
\[
P_B(A)=P(A\mid B)=\frac{P(A\cap B)}{P(B)}.
\]
\end{definition}

\begin{propriete}[Probabilité d'une intersection]
Si $P(B)\neq 0$ : $\;P(A\cap B)=P(B)\,P_B(A)$. De même, si $P(A)\neq 0$,
$P(A\cap B)=P(A)\,P_A(B)$.
\end{propriete}

\begin{definition}[Événements indépendants]
Deux événements $A$ et $B$ sont \textbf{indépendants} lorsque
\[
P(A\cap B)=P(A)\times P(B),
\]
ce qui équivaut, si $P(B)\neq 0$, à $P_B(A)=P(A)$ : la réalisation de $B$ ne
modifie pas la probabilité de $A$.
\end{definition}

\begin{remarque}
Ne pas confondre \textbf{incompatibles} ($A\cap B=\varnothing$) et
\textbf{indépendants} ($P(A\cap B)=P(A)P(B)$). Deux événements incompatibles de
probabilités non nulles ne sont \emph{jamais} indépendants.
\end{remarque}

\begin{exemple}
On lance deux dés équilibrés. $A$ : « le premier dé donne $6$ », $B$ : « la somme
vaut $7$ ». On a $P(A)=\frac16$, $P(B)=\frac{6}{36}=\frac16$ et $A\cap B$ : « $(6,1)$ »
soit $P(A\cap B)=\frac1{36}=P(A)P(B)$ : $A$ et $B$ sont \textbf{indépendants}.
\end{exemple}

% ═══════════════════════════════════════════════════════════════
\section{Probabilités totales et arbre pondéré}

\begin{definition}[Système complet d'événements]
Des événements $A_1,A_2,\dots,A_n$ de probabilités non nulles forment un
\textbf{système complet d'événements} lorsqu'ils sont deux à deux incompatibles
et que leur réunion est $\Omega$ : ils « partagent » $\Omega$ sans recouvrement.
\end{definition}

\begin{theoreme}[Formule des probabilités totales]
Si $A_1,\dots,A_n$ forment un système complet d'événements, alors pour tout
événement $B$ :
\[
P(B)=\sum_{i=1}^{n}P(A_i\cap B)=\sum_{i=1}^{n}P(A_i)\,P_{A_i}(B).
\]
\end{theoreme}

\begin{aretenir}
Un \textbf{arbre pondéré} traduit la formule des probabilités totales. Règles :
\begin{itemize}
\item la somme des probabilités des branches \emph{issues d'un même n\oe ud} vaut $1$ ;
\item la probabilité d'un \emph{chemin} est le \textbf{produit} des probabilités
rencontrées ;
\item la probabilité d'un événement est la \textbf{somme} des probabilités des
chemins qui le réalisent.
\end{itemize}
\end{aretenir}

\begin{illus}
\begin{tikzpicture}[>=Stealth,grow=right,level distance=26mm,
  every node/.style={font=\small},
  nd/.style={circle,draw=onbo,fill=onbsoft,inner sep=1.6pt},
  lf/.style={rectangle,draw=onbline,fill=white,inner sep=2.5pt}]
\node[nd] (r) {}
  child[level distance=30mm]{node[nd] (A) {$A$}
    child{node[lf] {$B$}    edge from parent node[above,onbgreen]{\scriptsize $0{,}9$}}
    child{node[lf] {$\overline B$} edge from parent node[below,onbmuted]{\scriptsize $0{,}1$}}
    edge from parent node[above,onbo]{\scriptsize $0{,}4$}}
  child[level distance=30mm]{node[nd] (Ab) {$\overline A$}
    child{node[lf] {$B$}    edge from parent node[above,onbgreen]{\scriptsize $0{,}3$}}
    child{node[lf] {$\overline B$} edge from parent node[below,onbmuted]{\scriptsize $0{,}7$}}
    edge from parent node[below,onbo]{\scriptsize $0{,}6$}};
\node[align=left,text width=4.3cm,anchor=west,onbink] at (5.4,0)
{\small $P(B)=0{,}4\times0{,}9+0{,}6\times0{,}3$ \\[2pt] $\phantom{P(B)}=0{,}36+0{,}18=0{,}54.$};
\end{tikzpicture}
\fig{Figure — Arbre pondéré : système complet $\{A,\overline A\}$ et probabilités totales.}
\end{illus}

\begin{exemple}
Une usine de Douala possède deux chaînes. La chaîne $A$ produit $40\%$ des pièces
avec $90\%$ de bonnes ($P_A(B)=0{,}9$), la chaîne $\overline A$ produit les
$60\%$ restants avec $30\%$ de bonnes. La probabilité qu'une pièce prise au
hasard soit bonne est, d'après l'arbre,
$P(B)=0{,}4\cdot0{,}9+0{,}6\cdot0{,}3=0{,}54$.
\end{exemple}

\begin{theoreme}[Probabilité des causes (formule de Bayes)]
Si $A_1,\dots,A_n$ est un système complet et $P(B)\neq 0$, alors pour tout $i$ :
\[
P_B(A_i)=\frac{P(A_i)\,P_{A_i}(B)}{\displaystyle\sum_{j=1}^{n}P(A_j)\,P_{A_j}(B)}.
\]
\end{theoreme}

% ═══════════════════════════════════════════════════════════════
\section{Variables aléatoires}

\begin{definition}[Variable aléatoire et loi de probabilité]
Une \textbf{variable aléatoire (v.a.)} $X$ est une application de $\Omega$ dans
$\R$ qui associe un réel à chaque issue. L'ensemble de ses valeurs est
$\{x_1,\dots,x_n\}$. La \textbf{loi de probabilité} de $X$ est la donnée des
probabilités $p_i=P(X=x_i)$, avec
\[
p_i\ge 0 \quad\text{et}\quad \sum_{i=1}^{n}p_i=1.
\]
\end{definition}

\begin{definition}[Fonction de répartition]
La \textbf{fonction de répartition} de $X$ est la fonction $F$ définie sur $\R$ par
\[
F(x)=P(X\le x).
\]
C'est une fonction \textbf{en escalier}, croissante, qui passe de $0$ à $1$ ;
elle augmente d'un « saut » de hauteur $p_i$ à chaque valeur $x_i$.
\end{definition}

\begin{exemple}
Soit $X$ le nombre de « piles » obtenus en lançant deux fois une pièce
équilibrée. Les issues équiprobables sont PP, PF, FP, FF, donc $P(X{=}0)=\frac14$,
$P(X{=}1)=\frac12$, $P(X{=}2)=\frac14$. La fonction de répartition vaut $0$ pour
$x<0$, $\frac14$ pour $0\le x<1$, $\frac34$ pour $1\le x<2$, et $1$ pour $x\ge 2$.
\end{exemple}

% ═══════════════════════════════════════════════════════════════
\section{Espérance, variance et écart-type}

\begin{definition}[Espérance, variance, écart-type]
Soit $X$ une v.a. de loi $P(X=x_i)=p_i$. On définit :
\[
E(X)=\sum_{i=1}^{n}p_i\,x_i,\qquad
V(X)=\sum_{i=1}^{n}p_i\big(x_i-E(X)\big)^2,\qquad
\sigma(X)=\sqrt{V(X)}.
\]
$E(X)$ est la \textbf{valeur moyenne} attendue, $V(X)\ge0$ mesure la
\textbf{dispersion} et $\sigma(X)$ l'\textbf{écart-type}.
\end{definition}

\begin{propriete}[Formule de König--Huygens et linéarité]
La variance se calcule plus rapidement par
\[
V(X)=E(X^2)-\big(E(X)\big)^2,\qquad\text{où}\quad E(X^2)=\sum_{i=1}^{n}p_i\,x_i^{2}.
\]
De plus, pour tous réels $a,b$ : $\;E(aX+b)=a\,E(X)+b$ et $V(aX+b)=a^{2}V(X)$.
\end{propriete}

\begin{exemple}
Pour le $X$ ci-dessus (deux lancers) :
$E(X)=0\cdot\frac14+1\cdot\frac12+2\cdot\frac14=1$, $\;E(X^2)=0\cdot\frac14+1\cdot\frac12+4\cdot\frac14=\frac32$,
donc $V(X)=\frac32-1^2=\frac12$ et $\sigma(X)=\frac{1}{\sqrt2}\approx0{,}71$.
\end{exemple}

% ═══════════════════════════════════════════════════════════════
\section{Schéma de Bernoulli et loi binomiale}

\begin{definition}[Épreuve et loi de Bernoulli]
Une \textbf{épreuve de Bernoulli} est une expérience à deux issues : le
\textbf{succès} $S$ (probabilité $p$) et l'\textbf{échec} $\overline S$
(probabilité $q=1-p$). La v.a. $X$ qui vaut $1$ en cas de succès et $0$ sinon
suit la \textbf{loi de Bernoulli} de paramètre $p$ : $E(X)=p$, $V(X)=p(1-p)$.
\end{definition}

\begin{definition}[Schéma de Bernoulli]
Un \textbf{schéma de Bernoulli} est la répétition de $n$ épreuves de Bernoulli
\textbf{identiques et indépendantes}, de même paramètre $p$.
\end{definition}

\begin{theoreme}[Loi binomiale]
Dans un schéma de $n$ épreuves de Bernoulli de paramètre $p$, la v.a. $X$ égale
au \textbf{nombre de succès} suit la \textbf{loi binomiale} $\mathcal B(n,p)$.
Pour tout entier $k$ avec $0\le k\le n$ :
\[
P(X=k)=\binom{n}{k}\,p^{k}(1-p)^{\,n-k}.
\]
Son espérance et sa variance sont
\[
E(X)=np,\qquad V(X)=np(1-p),\qquad \sigma(X)=\sqrt{np(1-p)}.
\]
\end{theoreme}

\begin{remarque}
Le coefficient $\binom nk$ compte les façons de placer les $k$ succès parmi les
$n$ épreuves ; $p^k(1-p)^{n-k}$ est la probabilité d'un chemin précis (par
indépendance). La loi binomiale combine ainsi dénombrement et indépendance.
\end{remarque}

\begin{illus}
\begin{tikzpicture}[>=Stealth]
\begin{axis}[width=9.4cm,height=5.4cm,ybar,bar width=12pt,
  xlabel={\small $k$ (nombre de succès)},ylabel={\small $P(X=k)$},
  xmin=-0.6,xmax=6.6,ymin=0,ymax=0.32,
  xtick={0,1,2,3,4,5,6},ytick={0,0.1,0.2,0.3},
  axis lines=left,enlarge x limits=0.05,
  every axis plot/.append style={fill=onbo,fill opacity=0.85,draw=onbo}]
\addplot coordinates {(0,0.0467)(1,0.1866)(2,0.3110)(3,0.2765)(4,0.1382)(5,0.0369)(6,0.0041)};
\end{axis}
\node[align=left,text width=3.4cm,anchor=west,onbink] at (9.5,1.7)
{\small Loi $\mathcal B(6\,;\,0{,}4)$ : $E(X)=2{,}4$. La barre la plus haute est en $k=2$.};
\end{tikzpicture}
\fig{Figure — Diagramme en bâtons de la loi binomiale $\mathcal B(6\,;\,0{,}4)$.}
\end{illus}

\begin{exemple}
Un QCM comporte $6$ questions à $4$ choix ; un candidat répond au hasard. Le
nombre $X$ de bonnes réponses suit $\mathcal B(6\,;\,0{,}25)$. La probabilité
d'avoir exactement $2$ bonnes réponses est
$P(X{=}2)=\binom62(0{,}25)^2(0{,}75)^4=15\times0{,}0625\times0{,}3164\approx0{,}297$.
On attend en moyenne $E(X)=6\times0{,}25=1{,}5$ bonne réponse.
\end{exemple}

% ═══════════════════════════════════════════════════════════════
\section{Méthodes \& savoir-faire}

\methode{Calculer une probabilité par dénombrement}
Décrire l'univers, choisir le bon modèle (liste, arrangement ou combinaison),
compter les cas favorables et les cas possibles, puis appliquer
$P(A)=\frac{\Card A}{\Card\Omega}$.
\begin{exemple}
Une urne contient $5$ boules rouges et $3$ boules vertes. On tire simultanément
$3$ boules. Probabilité d'obtenir exactement $2$ rouges :
\[
P=\frac{\binom52\binom31}{\binom83}=\frac{10\times3}{56}=\frac{30}{56}=\frac{15}{28}\approx0{,}54.
\]
\end{exemple}

\methode{Utiliser un arbre pondéré et les probabilités totales}
Construire l'arbre du système complet, multiplier le long des chemins, additionner
les chemins favorables. Pour une probabilité conditionnelle « renversée »,
appliquer Bayes.
\begin{exemple}
Reprenons l'usine ($P(A){=}0{,}4$, $P_A(B){=}0{,}9$, $P_{\overline A}(B){=}0{,}3$).
Une pièce tirée est bonne ; probabilité qu'elle vienne de $A$ :
\[
P_B(A)=\frac{P(A)P_A(B)}{P(B)}=\frac{0{,}4\times0{,}9}{0{,}54}=\frac{0{,}36}{0{,}54}=\frac23\approx0{,}67.
\]
\end{exemple}

\methode{Établir la loi de probabilité et la fonction de répartition d'une v.a.}
Déterminer les valeurs prises par $X$, calculer chaque $P(X=x_i)$, vérifier que
la somme vaut $1$, puis cumuler les probabilités pour obtenir $F(x)=P(X\le x)$.
\begin{exemple}
On tire $2$ boules simultanément d'une urne de $3$ rouges et $2$ noires ; $X$ =
nombre de rouges. $P(X{=}0)=\frac{\binom22}{\binom52}=\frac1{10}$,
$P(X{=}1)=\frac{\binom31\binom21}{10}=\frac6{10}$,
$P(X{=}2)=\frac{\binom32}{10}=\frac3{10}$ (somme $=1$). Donc $F(x)=0$ si $x<0$,
$\frac1{10}$ si $0\le x<1$, $\frac7{10}$ si $1\le x<2$, $1$ si $x\ge2$.
\end{exemple}

\methode{Calculer espérance, variance et écart-type}
Utiliser $E(X)=\sum p_i x_i$, puis $V(X)=E(X^2)-E(X)^2$ et $\sigma=\sqrt{V}$.
\begin{exemple}
Avec la loi précédente : $E(X)=0\cdot\frac1{10}+1\cdot\frac6{10}+2\cdot\frac3{10}=\frac{12}{10}=1{,}2$ ;
$E(X^2)=0+1\cdot\frac6{10}+4\cdot\frac3{10}=\frac{18}{10}=1{,}8$ ;
$V(X)=1{,}8-1{,}44=0{,}36$, $\sigma(X)=0{,}6$.
\end{exemple}

\methode{Reconnaître et utiliser la loi binomiale}
Vérifier les trois conditions : épreuves \emph{identiques}, \emph{indépendantes},
à \emph{deux issues}. Identifier $n$ et $p$, puis appliquer
$P(X{=}k)=\binom nk p^k(1-p)^{n-k}$, $E=np$, $V=np(1-p)$.
\begin{exemple}
On lance $5$ fois un dé ; succès = « obtenir un $6$ », $p=\frac16$. Le nombre $X$
de $6$ suit $\mathcal B(5\,;\,\frac16)$. Probabilité d'au moins un $6$ :
\[
P(X\ge1)=1-P(X{=}0)=1-\Big(\tfrac56\Big)^5=1-\tfrac{3125}{7776}\approx0{,}598.
\]
\end{exemple}

% ═══════════════════════════════════════════════════════════════
\section{Exercices}

\rubrique{Équiprobabilité et dénombrement}
\exo{}\diff{1} On tire une carte d'un jeu de $52$ cartes. Calculer la probabilité
d'obtenir : a) un roi ; b) un c\oe ur ; c) un roi ou un c\oe ur.

\exo{}\diff{1} Un sac contient $4$ jetons numérotés de $1$ à $4$. On en tire deux
simultanément. Quelle est la probabilité que la somme des numéros soit paire ?

\exo{}\diff{2} Une urne contient $6$ boules rouges et $4$ boules blanches. On tire
simultanément $3$ boules. Calculer la probabilité d'obtenir : a) trois rouges ;
b) exactement deux rouges ; c) au moins une blanche.

\rubrique{Probabilité conditionnelle, arbres et probabilités totales}
\exo{}\diff{2} Dans une classe de Terminale C, $60\%$ des élèves sont des filles.
$70\%$ des filles et $50\%$ des garçons sont internes. On choisit un élève au
hasard. a) Faire un arbre pondéré. b) Calculer la probabilité qu'il soit interne.
c) Sachant qu'il est interne, quelle est la probabilité que ce soit une fille ?

\exo{}\diff{2} Un test de dépistage détecte une maladie touchant $2\%$ d'une
population. Il est positif chez $95\%$ des malades et chez $4\%$ des non-malades.
Une personne est testée positive : quelle est la probabilité qu'elle soit
réellement malade ?

\exo{}\diff{3} Deux urnes : $U_1$ contient $3$ rouges et $2$ vertes, $U_2$
contient $1$ rouge et $4$ vertes. On lance un dé équilibré : s'il donne $6$, on
tire dans $U_1$, sinon dans $U_2$. a) Probabilité de tirer une boule rouge.
b) On a tiré une rouge ; probabilité qu'elle vienne de $U_1$.

\rubrique{Variables aléatoires, espérance et variance}
\exo{}\diff{2} On lance deux dés équilibrés et on note $X$ la somme des points
obtenus. a) Donner la loi de probabilité de $X$. b) Calculer $E(X)$.

\exo{}\diff{2} Un jeu coûte $200$ FCFA. On tire une boule d'une urne contenant
$2$ boules gagnantes (gain $500$ FCFA) et $8$ perdantes (gain nul). Soit $G$ le
gain algébrique du joueur. a) Donner la loi de $G$. b) Le jeu est-il favorable au
joueur (calculer $E(G)$) ? c) Calculer $\sigma(G)$.

\exo{}\diff{3} Une v.a. $X$ prend les valeurs $-1,0,2$ avec
$P(X{=}-1)=\frac14$, $P(X{=}0)=a$, $P(X{=}2)=\frac12$. a) Déterminer $a$.
b) Calculer $E(X)$, $V(X)$, $\sigma(X)$. c) Donner la fonction de répartition $F$.

\rubrique{Loi binomiale}
\exo{}\diff{2} Une machine produit des pièces dont $5\%$ sont défectueuses. On
prélève $10$ pièces (production assez grande pour supposer l'indépendance). Soit
$X$ le nombre de pièces défectueuses. a) Justifier que $X$ suit une loi binomiale.
b) Calculer $P(X{=}0)$, $P(X{=}1)$ et $P(X\ge1)$. c) Donner $E(X)$.

\exo{}\diff{2} Un footballeur réussit chacun de ses penalties avec la probabilité
$0{,}8$, indépendamment les uns des autres. Il en tire $4$. Soit $X$ le nombre de
buts. a) Donner la loi de $X$. b) Probabilité qu'il marque exactement $3$ buts.
c) Probabilité qu'il marque au moins $3$ buts. d) Calculer $E(X)$ et $\sigma(X)$.

\rubrique{Problème type Baccalauréat}
\exo{}\diff{3} \textbf{Épreuve composée, v.a. et loi binomiale.}
Une urne contient $5$ boules rouges et $3$ boules noires, indiscernables au
toucher.

\emph{Partie A.} On tire simultanément $3$ boules de l'urne. Soit $X$ le nombre
de boules rouges tirées.
\begin{enumerate}[label=\arabic*)]
\item Déterminer la loi de probabilité de $X$.
\item Calculer $E(X)$ et $V(X)$.
\end{enumerate}

\emph{Partie B.} On remet les boules dans l'urne. On effectue maintenant $4$
tirages \textbf{successifs avec remise} d'une boule. Soit $Y$ le nombre de boules
rouges obtenues au cours de ces $4$ tirages.
\begin{enumerate}[label=\arabic*),start=3]
\item Justifier que $Y$ suit une loi binomiale dont on précisera les paramètres.
\item Calculer $P(Y{=}2)$ et $P(Y\ge1)$.
\item Déterminer $E(Y)$ et $\sigma(Y)$.
\end{enumerate}

% ═══════════════════════════════════════════════════════════════
\section{Corrigés}

\corrige{de l'exercice 1}
a) $P=\frac{4}{52}=\frac1{13}$.\;
b) $P=\frac{13}{52}=\frac14$.\;
c) « roi ou c\oe ur » : $P=\frac{4}{52}+\frac{13}{52}-\frac{1}{52}=\frac{16}{52}=\frac{4}{13}$ (le roi de c\oe ur est compté une fois).

\corrige{de l'exercice 2}
$\Card\Omega=\binom42=6$. Somme paire $\Leftrightarrow$ les deux numéros de même
parité : paires $\{1,3\}$ et $\{2,4\}$, soit $2$ cas. $P=\frac26=\frac13$.

\corrige{de l'exercice 3}
$\Card\Omega=\binom{10}{3}=120$.\;
a) $\frac{\binom63}{120}=\frac{20}{120}=\frac16$.\;
b) $\frac{\binom62\binom41}{120}=\frac{15\times4}{120}=\frac{60}{120}=\frac12$.\;
c) « au moins une blanche » = contraire de « trois rouges » : $1-\frac16=\frac56$.

\corrige{de l'exercice 4}
Notons $F$ (fille), $G$ (garçon), $I$ (interne). $P(F)=0{,}6$, $P(G)=0{,}4$,
$P_F(I)=0{,}7$, $P_G(I)=0{,}5$.\;
b) $P(I)=0{,}6\times0{,}7+0{,}4\times0{,}5=0{,}42+0{,}20=0{,}62$.\;
c) $P_I(F)=\dfrac{P(F)P_F(I)}{P(I)}=\dfrac{0{,}42}{0{,}62}=\dfrac{42}{62}=\dfrac{21}{31}\approx0{,}68$.

\corrige{de l'exercice 5}
$M$ : malade, $T$ : test positif. $P(M)=0{,}02$, $P_M(T)=0{,}95$,
$P_{\overline M}(T)=0{,}04$.
\[
P(T)=0{,}02\times0{,}95+0{,}98\times0{,}04=0{,}019+0{,}0392=0{,}0582.
\]
\[
P_T(M)=\frac{0{,}019}{0{,}0582}\approx0{,}326.
\]
Malgré un test « fiable », à peine un tiers des positifs sont réellement malades
(la maladie est rare) : c'est un piège classique des dépistages.

\corrige{de l'exercice 6}
Notons $R$ : « rouge ». $P(U_1)=\frac16$, $P(U_2)=\frac56$,
$P_{U_1}(R)=\frac35$, $P_{U_2}(R)=\frac15$.\;
a) $P(R)=\frac16\cdot\frac35+\frac56\cdot\frac15=\frac{3}{30}+\frac{5}{30}=\frac{8}{30}=\frac{4}{15}$.\;
b) $P_R(U_1)=\dfrac{\frac16\cdot\frac35}{\frac{4}{15}}=\dfrac{\frac{3}{30}}{\frac{8}{30}}=\dfrac{3}{8}$.

\corrige{de l'exercice 7}
a) $X$ prend les valeurs $2,3,\dots,12$ avec $\Card\Omega=36$. Les probabilités
(numérateur = nombre de couples) :
\[
\begin{array}{c|ccccccccccc}
k & 2&3&4&5&6&7&8&9&10&11&12\\\hline
P(X{=}k)&\frac1{36}&\frac2{36}&\frac3{36}&\frac4{36}&\frac5{36}&\frac6{36}&\frac5{36}&\frac4{36}&\frac3{36}&\frac2{36}&\frac1{36}
\end{array}
\]
b) Par symétrie autour de $7$ : $E(X)=7$.

\corrige{de l'exercice 8}
$G$ vaut $500-200=300$ (gain net en cas de succès) ou $0-200=-200$ (perte).
$P(G{=}300)=\frac{2}{10}=\frac15$, $P(G{=}-200)=\frac{8}{10}=\frac45$.\;
b) $E(G)=300\cdot\frac15+(-200)\cdot\frac45=60-160=-100$ FCFA : le jeu est
\textbf{défavorable} au joueur (perte moyenne de $100$ FCFA par partie).\;
c) $E(G^2)=300^2\cdot\frac15+200^2\cdot\frac45=18000+32000=50000$ ;
$V(G)=50000-(-100)^2=40000$ ; $\sigma(G)=200$ FCFA.

\corrige{de l'exercice 9}
a) $\frac14+a+\frac12=1\Rightarrow a=\frac14$.\;
b) $E(X)=-1\cdot\frac14+0\cdot\frac14+2\cdot\frac12=-\frac14+1=\frac34$.
$E(X^2)=1\cdot\frac14+0+4\cdot\frac12=\frac14+2=\frac94$.
$V(X)=\frac94-\frac9{16}=\frac{36-9}{16}=\frac{27}{16}$, $\sigma(X)=\frac{\sqrt{27}}{4}=\frac{3\sqrt3}{4}\approx1{,}30$.\;
c) $F(x)=0$ si $x<-1$ ; $\frac14$ si $-1\le x<0$ ; $\frac12$ si $0\le x<2$ ; $1$ si $x\ge2$.

\corrige{de l'exercice 10}
a) Chaque pièce est défectueuse ou non (deux issues), avec $p=0{,}05$, les $10$
prélèvements étant indépendants : $X\sim\mathcal B(10\,;\,0{,}05)$.\;
b) $P(X{=}0)=(0{,}95)^{10}\approx0{,}599$ ;
$P(X{=}1)=\binom{10}{1}(0{,}05)(0{,}95)^9=10\times0{,}05\times0{,}6302\approx0{,}315$ ;
$P(X\ge1)=1-(0{,}95)^{10}\approx0{,}401$.\;
c) $E(X)=10\times0{,}05=0{,}5$.

\corrige{de l'exercice 11}
a) $X\sim\mathcal B(4\,;\,0{,}8)$ : $P(X{=}k)=\binom4k(0{,}8)^k(0{,}2)^{4-k}$.\;
b) $P(X{=}3)=\binom43(0{,}8)^3(0{,}2)=4\times0{,}512\times0{,}2=0{,}4096$.\;
c) $P(X\ge3)=P(X{=}3)+P(X{=}4)=0{,}4096+(0{,}8)^4=0{,}4096+0{,}4096=0{,}8192$.\;
d) $E(X)=4\times0{,}8=3{,}2$ ; $V(X)=4\times0{,}8\times0{,}2=0{,}64$, $\sigma(X)=0{,}8$.

\corrige{du problème}
\textbf{Partie A.} $\Card\Omega=\binom83=56$. $X$ (nombre de rouges parmi $3$
tirées) prend les valeurs $0,1,2,3$ :
\[
P(X{=}0)=\tfrac{\binom33}{56}=\tfrac1{56},\quad
P(X{=}1)=\tfrac{\binom51\binom32}{56}=\tfrac{15}{56},\quad
P(X{=}2)=\tfrac{\binom52\binom31}{56}=\tfrac{30}{56},\quad
P(X{=}3)=\tfrac{\binom53}{56}=\tfrac{10}{56}.
\]
(Somme : $1+15+30+10=56$, donc $1$.)\;
2) $E(X)=\dfrac{0+15+60+30}{56}=\dfrac{105}{56}=\dfrac{15}{8}=1{,}875$.
$E(X^2)=\dfrac{0+15+120+90}{56}=\dfrac{225}{56}$, donc
$V(X)=\dfrac{225}{56}-\Big(\dfrac{15}{8}\Big)^2=\dfrac{225}{56}-\dfrac{225}{64}
=225\Big(\dfrac{64-56}{3584}\Big)=\dfrac{1800}{3584}=\dfrac{225}{448}\approx0{,}502$.

\textbf{Partie B.} 3) À chaque tirage avec remise, la composition est inchangée :
les $4$ tirages sont identiques et indépendants, à deux issues (rouge / non
rouge) avec $p=\frac58$. Donc $Y\sim\mathcal B\big(4\,;\,\frac58\big)$.\;
4) $P(Y{=}2)=\binom42\Big(\frac58\Big)^2\Big(\frac38\Big)^2
=6\times\frac{25}{64}\times\frac{9}{64}=\frac{1350}{4096}\approx0{,}330$.
$P(Y\ge1)=1-P(Y{=}0)=1-\Big(\frac38\Big)^4=1-\frac{81}{4096}=\frac{4015}{4096}\approx0{,}980$.\;
5) $E(Y)=4\times\frac58=\frac52=2{,}5$ ;
$V(Y)=4\times\frac58\times\frac38=\frac{60}{64}=\frac{15}{16}$,
$\sigma(Y)=\frac{\sqrt{15}}{4}\approx0{,}968$.
